\begin{table}[htbp]
  \centering
  \caption{Empirical assessment of the three thesis hypotheses}
  \label{tab:hypothesis-assessment}
  \small
  \setlength{\tabcolsep}{5pt}
  \begin{tabularx}{\textwidth}{@{}>{\raggedright\arraybackslash}p{0.24\textwidth}>{\raggedright\arraybackslash}X>{\raggedright\arraybackslash}p{0.20\textwidth}@{}}
    \toprule
    Hypothesis & Evidence & Assessment \\
    \midrule
    H1: Housing-credit persistence & House-purchase lending growth is positive on impact and peaks at h6. Its accumulated response excludes zero from h0 through h12 under the HAC interval and both bootstraps. & Supported at short and medium horizons. \\
    H2: Compensation-linked attenuation & The canonical wage tracker moves at short and medium horizons but loses h24 cumulative exclusion, while employment expectations reverse. The German industry wage bill, a sectoral robustness proxy, keeps its sign and excludes zero from h1 through h24 under every interval method, which is the proxy boundary on this verdict. & Supported with proxy boundary. \\
    H3: Differential financial-intermediation transmission & Housing-credit growth and lending spreads are stronger and more persistent than NFC credit quantities and the euro-area wage tracker. & Conditionally supported. \\
    \bottomrule
  \end{tabularx}
  \vspace{0.25em}
  {\footnotesize Notes: H1 to H3 are Hypotheses 1 to 3 of \Cref{chap:introduction}, derived in \Cref{sec:hypothesis-derivation}. CI refers to the HAC interval under the independence approximation, bootstrap to the joint moving-block and wild intervals in \Cref{tab:bootstrap-cumulative}. Source: own assessment based on own calculations.}
\end{table}
